% Backup slides, after "Thank you!": for questions, and for the specialists
% in the room. Numbered B1, B2, ... (preamble: \backuppart), not counted in
% the talk's total. "Quicksort, Step by Step" and "Knapsack, Step by Step" are
% shared with the seminar deck: keep them identical (../check-hylos-sync.py,
% in tex/).
\backuppart

\begin{frame}
  \frametitle{Hylomorphisms: Generality, Laws, Parallelism}
  \small
  \begin{itemize}
    \item \textbf{General.} Conjugate hylomorphisms (adjoint folds and
      unfolds) cover paramorphisms (primitive recursion), histomorphisms
      (course-of-values recursion, as in dynamic programming),
      accumulations and mutual recursion. The solution is unique when the
      coalgebra is \emph{recursive}, e.g.\ when it terminates.
    \item \textbf{Laws.} Fusion and deforestation remove passes and
      intermediate structures. GHC implements shortcut fusion
      (\code{foldr}/\code{build}) as rewrite rules; Takano and Meijer
      generalise it to hylomorphisms (acid rain); the HYLO system finds
      hylomorphisms in ordinary recursive definitions, then fuses them.
    \item \textbf{Parallelism.} The sub-problems of a layer are
      independent: parallel skeletons are hylomorphisms (ICFP 2016), and
      so are their cost models (FGCS 2018).
  \end{itemize}
  \slideref{P. Wadler. \textbf{Deforestation: Transforming Programs to Eliminate Trees}. TCS 1990. \rfsep
    A. Gill, J. Launchbury, S. Peyton Jones. \textbf{A Short Cut to Deforestation}. FPCA 1993. \rfsep
    A. Takano, E. Meijer. \textbf{Shortcut Deforestation in Calculational Form}. FPCA 1995. \rfsep
    Z. Hu, H. Iwasaki, M. Takeichi. \textbf{Deriving Structural Hylomorphisms from Recursive Definitions}. ICFP 1996. \rfsep
    Y. Onoue, Z. Hu, H. Iwasaki, M. Takeichi. \textbf{A Calculational Fusion System HYLO}. Algorithmic Languages and Calculi, 1997. \rfsep
    V. Capretta, T. Uustalu, V. Vene. \textbf{Recursive Coalgebras from Comonads}. Inf. Comput. 2006. \rfsep
    R. Hinze, N. Wu, J. Gibbons. \textbf{Conjugate Hylomorphisms, or: The Mother of All Structured Recursion Schemes}. POPL 2015.}
  \speakernote{Conjugate hylos: adjunctions give para, histo, accumulations,
    mutual recursion \\
    Uniqueness: recursive coalgebras (Part II) \\
    GHC RULES: foldr/build; acid rain; HYLO system \\
    Skeletons and cost models: ICFP 2016, FGCS 2018}
\end{frame}

\begin{frame}[vertical-alignment = top]
  \frametitle{The Third Homomorphism Theorem, and After}
  \vspace*{2mm}
  \begin{center}
    \scriptsize
    \renewcommand\arraystretch{1.05}
    \begin{tabular}{@{}l@{\quad}l@{\quad}>{\raggedright\arraybackslash}p{72mm}@{}}
      1987 & Bird & theory of lists: 1st and 2nd theorems (map, then
        reduce; leftward and rightward reductions) \\
      1989 & Meertens & 3rd theorem proved, as conjectured by Bird \\
      1996 & Gibbons, JFP & 3rd theorem published; mergesort from
        insertion sort \\
      1996 & Gorlatch, Euro-Par & scans and homomorphisms, parallelised
        systematically \\
      1997 & Hu, Iwasaki, Takeichi, TOPLAS & homomorphisms by tupling
        with extra state (lifting) \\
      2007 & Morita et al., PLDI & combiners derived automatically, by
        weak right inverses \\
      2009 & Morihata et al., POPL & the third theorem on trees (downward
        and upward) \\
      2011 & Mu, Morihata, ICFP & generalised and dualised \\
      2011 & Chi, Mu, APLAS & combiners constructed from a proof that the
        two loops agree \\
      2017 & Farzan, Nicolet, PLDI & combiners and lifting synthesised for
        loops \\
      2019 & Farzan, Nicolet, PLDI & nested loops, one level at a time \\
      2021 & Farzan, Nicolet, PLDI & phased synthesis of divide and
        conquer \\
    \end{tabular}
  \end{center}
  \speakernote{From pen-and-paper theorems to tools \\
    Lifting = extra state, as in the Java example (tupling) \\
    Morita et al.: the combiner by program inversion \\
    Chi, Mu: from the proof that the loops agree, as in Part I \\
    Farzan and Nicolet: synthesis (ParSynt); nested loops in Part III}
\end{frame}

\begin{frame*}[vertical-alignment = top]
  \frametitle{ITP 2025's Extracted Quicksort, and Its Closures}
  \vspace*{3mm}
  \ocamlblock[before upper = {\fvset{fontsize = \scriptsize}}]{itp-qsort.ml}
  \footnotesize
  \begin{itemize}
    \item Type-safe, no \code{Obj.magic}: one type of positions
      (\code{Lbranch}, \code{Rbranch}) for every shape.
    \item A \textbf{closure} per call (\code{x0}): \code{hylo} takes apart
      the layer \code{c x} with a \code{match} that sits \emph{outside} the
      coalgebra's own \code{match} on the input. Extraction does not push
      it inside, so the children stay a function.
    \item The fix: \textbf{CPS coalgebras} (\code{CCoalg}). The coalgebra
      passes its layer to a continuation inside each of its branches, with
      a proof that it is \code{c x} (for termination). The algebra meets a
      known layer, and \code{x0 Lbranch} becomes \code{qsort\_double
      l0}; \code{hylo\_cps\_hylo} proves the two equal.
  \end{itemize}
  \rocqtag
  \speakernote{Real output of the paper's version (Rocq 9 port), renamed \\
    x0: the children as a function, a closure each call \\
    Cause: a match on a match; extraction does not commute them \\
    CPS: the algebra runs inside the coalgebra's branches; cps\_from\_coalg
    builds the CPS version from c}
\end{frame*}

\begin{frame*}[vertical-alignment = top]
  \frametitle{Quicksort, Step by Step}
  \vspace*{4mm}
  \begin{onlyenv}<1>
    \fvset{lastline = 1}%
    \rocqblock{qsort-steps.v}
  \end{onlyenv}
  \begin{onlyenv}<2>
    \fvset{lastline = 3}%
    \rocqblock{qsort-steps.v}
  \end{onlyenv}
  \begin{onlyenv}<3>
    \fvset{lastline = 5}%
    \rocqblock{qsort-steps.v}
  \end{onlyenv}
  \begin{onlyenv}<4->
    \fvset{lastline = 7}%
    \rocqblock{qsort-steps.v}
  \end{onlyenv}
  \par
  \speakernote{Each step: one law, the term Rocq shows after it \\
    map\_merge\_fusion: doubling moves inside merge \\
    hylo\_ana, deforest: the tree between unfold and fold disappears (one step) \\
    hylo\_cps\_hylo: continuation-passing coalgebra, so the code has no closures}
\end{frame*}

\begin{frame*}[vertical-alignment = top]
  \frametitle{Knapsack, Step by Step}
  \vspace*{4mm}
  \begin{onlyenv}<1>
    \fvset{lastline = 1}%
    \rocqblock{knap-steps.v}
  \end{onlyenv}
  \begin{onlyenv}<2>
    \fvset{lastline = 3}%
    \rocqblock{knap-steps.v}
  \end{onlyenv}
  \begin{onlyenv}<3>
    \fvset{lastline = 5}%
    \rocqblock{knap-steps.v}
  \end{onlyenv}
  \begin{onlyenv}<4->
    \fvset{lastline = 10}%
    \rocqblock{knap-steps.v}
  \end{onlyenv}
  \par
  \speakernote{The goal after each rewrite, as Rocq prints it \\
    reenumerate: from items to capacities, lookups jump into the array, O(1) \\
    dyna\_loop: the table is built by a loop counting up \\
    Last line: a loop over the capacities, one step each}
\end{frame*}

\begin{frame}
  \frametitle{The Laws of the Knapsack Calculation}
  \small
  \begin{itemize}
    \item \code{reenumerate}: a \code{hylo} becomes a
      \textbf{dynamorphism} (\code{dyna}) over a chain \code{e} (here
      \code{out\_int\_c}: $c, c - 1, \ldots, 0$): every result is computed
      once, kept in a table, and looked up. Side condition (\code{Hlook}):
      the algebra gets its children's results by looking back along
      \code{e}.
    \item \code{look\_with (w := fast\_drop) fast\_drop\_walk}: looking up
      $K(c - w)$ walks back $w$ entries; \code{fast\_drop} jumps there in
      $O(1)$ (an array), proved equal to the walk.
    \item \code{dyna\_loop}: a dynamorphism over a count-down is a
      \code{loop} counting up, one \code{memo\_step} per capacity.
    \item \code{dyna\_iso}: the result does not depend on how the table is
      represented (linked, or arrays).
  \end{itemize}
  \slideref{J. Kabanov, V. Vene. \textbf{Recursion Schemes for Dynamic Programming}. MPC 2006. \rfsep
    R. Hinze, N. Wu, J. Gibbons. \textbf{Conjugate Hylomorphisms, or: The Mother of All Structured Recursion Schemes}. POPL 2015.}
  \rocqtag
  \speakernote{Dynamorphism: a hylo with a memo table (Kabanov, Vene) \\
    reenumerate: Hlook, the children are found along the chain \\
    fast\_drop: O(1) jump, proved equal to the walk \\
    dyna\_loop: bottom-up loop; dyna\_iso: any table}
\end{frame}

\begin{frame*}
  \frametitle{Associative Combiners, for Any Data Type}
  \rocqblock{malg.v}
  \rocqblock{malg-fusion.v}
  \small
  \begin{itemize}
    \item \code{malg}: a \textbf{monoid algebra}. It combines a node's own
      value (\code{leaf s}, optional) with its children's results, in the
      order \code{members F s}, with the monoid \code{M}.
    \item List homomorphisms, for \textbf{any} data type and
      \textbf{any} divide step. \code{malg\_fusion}: any monoid
      homomorphism fuses; its side condition, \code{malg\_hom}, is proved
      once, for every shape.
    \item Instances: Bird's map and reduce promotion; maximum segment sum,
      from cubic (\code{maximum . map sum . segs}) to linear (Horner's
      rule, a scan, tupling), by \code{Derive}.
  \end{itemize}
  \rocqtag
  \speakernote{malg = monoid algebra: combine a node's value with its
    children's results, in order \\
    malg\_fusion: one-line proof, from malg\_hom \\
    Part I's homomorphisms, generalised \\
    Maximum segment sum: Bird's calculation, machine-checked (examples/MSS.v)}
\end{frame*}

\begin{frame}
  \frametitle{What Each Approach Trusts: ICFP 2016, Gibbon, ParSynt, Rocq}
  \begin{center}
    \small
    \renewcommand\arraystretch{1.35}
    \begin{tabular}{@{}llll@{}}
      & \textbf{the laws} & \textbf{each program} & \textbf{trusted} \\
      \hline
      ICFP 2016 & axioms & type checking & laws, code generator \\
      Gibbon & compiler passes & --- & the compiler \\
      ParSynt & paper proofs & join: bounded check, Dafny &
        translation, code generator \\
      \textbf{Rocq} & \textbf{proved once} & \textbf{every step checked} &
        \textbf{extraction or printer} \\
    \end{tabular}
  \end{center}
  \begin{itemize}[<+->]
    \item Same laws, one small trusted base.
    \item Clean extraction: the code that \textbf{runs} is a direct,
      readable transcription of what was \textbf{proved}.
  \end{itemize}
  \ongoingtag
  \speakernote{If asked how the leads compare on trust \\
    ICFP'16: laws assumed \\
    Gibbon: correct by careful engineering, unverified \\
    ParSynt: theorems on paper; each join checked; loop translation not
    verified \\
    Ours: laws are Rocq theorems; derivation is a certificate \\
    Without clean extraction, this is only a proof about a model}
\end{frame}

\begin{frame*}[vertical-alignment = top]
  \frametitle{Packed Traversals: End Witnesses, Offsets, Layout}
  \vspace*{4mm}
  \mainblock[before upper = {\fvset{fontsize = \scriptsize}}]{packed-sum.c}
  \scriptsize
  \begin{itemize}
    \item \textbf{End witnesses}: \code{sum} returns where its subtree
      ends, so the caller knows where the next one starts; LoCal makes
      these locations part of the types.
    \item \textbf{Random access}: to reach the right subtree without
      reading the left one, the buffer can also store offsets, at some cost
      in space; parallel traversals rely on them.
    \item \textbf{Layout}: which field comes first is a choice; Marmoset
      picks it with a cost model.
  \end{itemize}
  \slideref{M. Vollmer et al. \textbf{Compiling Tree Transforms to Operate on Packed Representations}. ECOOP 2017. \rfsep
    M. Vollmer et al. \textbf{LoCal: A Language for Programs Operating on Serialized Data}. PLDI 2019. \rfsep
    C. Koparkar et al. \textbf{Efficient Tree-Traversals: Reconciling Parallelism and Dense Data Representations}. ICFP 2021. \rfsep
    V. Singhal et al. \textbf{Optimizing Layout of Recursive Datatypes with Marmoset}. ECOOP 2024.}
  \ongoingtag
  \speakernote{The sum on the buffer (hand-written, not Gibbon's output) \\
    End witnesses: the cursor returned by each call \\
    Offsets: random access, needed for parallel traversals \\
    Order of traversal = order of layout}
\end{frame*}

